\documentclass[10pt,a4paper]{extarticle}
\pagenumbering{gobble}
\usepackage[left=0.62in,top=0.45in,right=0.62in,bottom=0.45in]{geometry}
\usepackage[T1]{fontenc}
\usepackage[utf8]{inputenc}
\usepackage{helvet}
\renewcommand{\familydefault}{\sfdefault}
\usepackage[turkish,shorthands=off]{babel}
\usepackage{xcolor}
\usepackage{titlesec}
\usepackage{enumitem}
\usepackage{array}
\usepackage{graphicx}
\usepackage[unicode]{hyperref}

\definecolor{accent}{HTML}{1A365D}
\definecolor{muted}{HTML}{4A5568}
\definecolor{rulecol}{HTML}{B8C4D4}

\hypersetup{colorlinks=true,urlcolor=accent,linkcolor=accent,
            pdftitle={Enes Orhan - Mobile and Backend Engineer},pdfauthor={Enes Orhan}}

%% ---- Bölüm hiyerarşisi -------------------------------------------------
\titleformat{\section}
  {\large\bfseries\scshape\color{accent}\raggedright}{}{0em}{}
  [\vspace{-1pt}{\color{rulecol}\titlerule[0.9pt]}]
\titlespacing*{\section}{0pt}{10pt}{4pt}

%% ---- Yapı taşları ------------------------------------------------------
\newcommand{\sep}{\;\textperiodcentered\;}
\newcommand{\ca}{\textasciitilde}

% Şirket satırı, konum sağa yaslı
\newcommand{\entry}[2]{%
\noindent\begin{tabular*}{\textwidth}{@{\extracolsep{\fill}}lr}
\textbf{#1} & #2\\
\end{tabular*}\vspace{1pt}
}
% Unvan satırı, tarihler sağa yaslı
\newcommand{\subentry}[2]{%
\noindent\begin{tabular*}{\textwidth}{@{\extracolsep{\fill}}lr}
\textit{\color{muted}#1} & \textit{\color{muted}#2}\\
\end{tabular*}\vspace{2pt}
}
% Proje altındaki teknoloji şeridi
\newcommand{\tech}[1]{{\footnotesize\color{muted}#1}}

\newlist{bullets}{itemize}{1}
\setlist[bullets]{leftmargin=1.1em,label=\textbullet,itemsep=2.5pt,topsep=2pt,parsep=0pt}

\begin{document}

%% ======================= BAŞLIK =======================
\begin{center}
{\fontsize{23}{27}\selectfont\bfseries\color{accent}Enes Orhan}\\[4pt]
{\large\color{muted}Mobile \& Backend Engineer}\\[8pt]
{\small
\href{tel:+905077561488}{+90 507 756 14 88}\sep
\href{mailto:enesorhann404@gmail.com}{enesorhann404@gmail.com}\sep
\href{https://linkedin.com/in/enesorhann}{linkedin.com/in/enesorhann}\sep
\href{https://github.com/enesorhann}{github.com/enesorhann}\sep
\href{https://enesorhan.com}{enesorhan.com}}
\end{center}

%% ======================= ÖZET =======================
\section{Özet}
Kullanıcı odaklı mobil uygulamalar ve ölçeklenebilir backend sistemleri geliştiriyorum. Mobil ürün
geliştirmeden backend servisleri ve altyapı yönetimine kadar uçtan uca geliştirme deneyimine sahibim.

%% ======================= DENEYİM =======================
\section{Deneyim}

\entry{Freelance Software Engineer}{Uzaktan}
\subentry{Software Engineer}{Ekim 2025 -- Şubat 2026}
\begin{bullets}
\item 3.000+ şirket ve \ca{}100 aktif kullanıcı için RBAC tabanlı çok kiracılı Django REST backend kurdum.
\item Param ve PayTR için işlem doğrulama ve callback yönetimi içeren ödeme altyapısı geliştirdim.
\item Postfix + Dovecot ile SPF, DKIM ve DMARC destekli production mail sunucusu yapılandırıp işlettim.
\item kubeadm ile Kubernetes cluster kurdum; dağıtım ve hata ayıklamayı kubectl ile yönettim.
\end{bullets}

\vspace{6pt}
\entry{Fiyator A.Ş. \textit{(Ankara, Türkiye)}\sep\href{https://fiyator.com}{fiyator.com}}{Uzaktan}
\subentry{Software Engineer Intern}{Haziran 2025 -- Ekim 2025}
\begin{bullets}
\item 14M+ ürün kaydı, 12M+ görsel ve 529K+ kategori içeren backend sistemlerinde çalıştım.
\item PostgreSQL indeksleme, Redis ve Elasticsearch ile API süresini \ca{}450ms'den \ca{}200ms'e düşürdüm.
\item Semantic Transformers ile \%95 doğrulukta Django + FastAPI kategori eşleştirme servisi geliştirdim.
\item Redis önbellekli FAISS vektör arama ve Ollama (gemma3n:e2b) çalıştıran FastAPI servisi tasarladım.
\item Flutter mobil yeniden tasarımını yürüttüm ve 45+ fonksiyonel/UI sorununu dokümante ettim.
\item Süresi dolmuş premium kullanıcıların erişimini koruyan abonelik doğrulama hatasını bulup çözdüm.
\end{bullets}

%% ======================= PROJELER =======================
\section{Projeler}
\begin{bullets}[itemsep=5pt]

\item \textbf{BidsForShips}\sep\href{https://bidsforships.com}{bidsforships.com}\\
Denizcilik sektörü için geliştirilmiş çok kiracılı B2B ihale ve teklif platformu.\\[2pt]
\tech{Django\sep PostgreSQL\sep Redis\sep Celery\sep Elasticsearch\sep Docker\sep Kubernetes\sep Hetzner}

\item \textbf{Zerbook}\hfill{\color{muted}Devam ediyor}\\
Flutter ile geliştirdiğim altın portföy takip uygulaması. .NET backend ve mağaza süreçlerini de yönetiyorum.\\[2pt]
\tech{Flutter\sep Dart\sep .NET\sep PostgreSQL\sep Firebase\sep Docker\sep Kubernetes\sep App Store\sep Google Play}

\item \textbf{Entegram}\sep\href{https://apps.apple.com/tr/app/entegram-kayd\%C4\%B1r-\%C3\%B6\%C4\%9Fren/id6766254758?l=tr}{App Store} / \href{https://play.google.com/store/apps/details?id=com.enesorhan.gibigram.gibigram_mobile}{Play Store}\\
Kısa ve kaydırmalı derslerden oluşan İngilizce öğrenme uygulaması; Flutter ve ASP.NET Core.\\[2pt]
\tech{Flutter\sep ASP.NET Core\sep PostgreSQL\sep Firebase\sep Docker\sep Kubernetes}

\item \textbf{Mail Sender Smart Edition}\sep\href{https://mailsendersmartedition.com}{mailsendersmartedition.com}\\
Yapay zekâ destekli toplu e-posta yönetim platformu.\\[2pt]
\tech{Flutter\sep Django\sep PostgreSQL\sep Gemini 2.5 Flash\sep Postfix + Dovecot\sep Docker\sep Kubernetes}

\item \textbf{CompAtlas}\sep\href{https://compatlas.com.tr/}{compatlas.com.tr}\hfill{\color{muted}Devam ediyor}\\
Türkiye'nin 81 ilindeki 17K+ teknoloji firmasını kapsayan, tek başıma geliştirdiğim veri platformu.\\[2pt]
\tech{Next.js\sep Supabase\sep SEO/GEO\sep OSM\sep Google Places\sep Overture Maps}

\item \textbf{Özyalçınlar Kereste}\sep React ile geliştirdiğim kurumsal kereste firması web sitesi {\color{muted}(yapım aşamasında)}.

\end{bullets}

%% ======================= YETENEKLER =======================
\section{Teknik Yetenekler}
\renewcommand{\arraystretch}{1.2}
\noindent\begin{tabular}{@{}>{\raggedright\arraybackslash}p{0.16\textwidth}@{}>{\raggedright\arraybackslash}p{0.84\textwidth}@{}}
\textbf{Diller} & Python, Dart, C\#, Kotlin, Java\\
\textbf{Mobil \& Web} & Flutter, Firebase, Next.js, React, Supabase, App Store \& Google Play yayını\\
\textbf{Backend} & Django, DRF, FastAPI, ASP.NET Core, Celery, REST API, RBAC, JWT\\
\textbf{Veritabanları} & PostgreSQL, Redis, Elasticsearch, SQLite\\
\textbf{DevOps} & Docker, Kubernetes (kubeadm), GitHub Actions, Linux, CI/CD, Hetzner\\
\textbf{AI/ML} & LLM entegrasyonu (Ollama, Gemini), FAISS vektör arama, Semantic Transformers\\
\end{tabular}

%% ======================= EĞİTİM =======================
\section{Eğitim \& Sertifikalar}
\noindent\textbf{Konya Teknik Üniversitesi}\sep Bilgisayar Mühendisliği (Lisans) {\color{muted}\sep Not Ort. 3.40 / 4.00}
\hfill {\color{muted}Ağu 2022 -- Haz 2026}\\[4pt]
\noindent\textbf{Sertifikalar}\quad
\href{https://drive.google.com/file/d/1btYql9fKCvBwqHeh06F_NE257FhceNCB/view}{Sertifika portföyü}\sep
\href{https://drive.google.com/file/d/12AcWLnSGEfcZUWUb945n6nD-iPB7aIA9/view?usp=sharing}{Teknofest Katılım Belgesi}

\end{document}
